\answersfor{ch:material}

\begin{solution}{ex:material:1}
Position of the 99.75th percentile:
\begin{align*}
N-1 &= 2{,}400{,}000-1 = 2{,}399{,}999,\\
\text{pos} &= 0.9975\times2{,}399{,}999\\
&= 2{,}399{,}999-0.0025\times2{,}399{,}999\\
&= 2{,}399{,}999-5{,}999.9975\\
&= 2{,}393{,}999.0025,\\
\lceil\text{pos}\rceil &= 2{,}394{,}000,\\
n_{\text{flag}} &= 2{,}400{,}000-2{,}394{,}000 = 6{,}000.
\end{align*}
That is $0.25\%$ of $2{,}400{,}000$, as intended.
\end{solution}

\begin{solution}{ex:material:2}
Rescale each map:
\begin{align*}
\tilde{s}_1 &= \frac{0.4-0.1}{1.6-0.1} = \frac{0.3}{1.5} = 0.20,\\
\tilde{s}_2 &= \frac{12-0}{60-0} = \frac{12}{60} = 0.20.
\end{align*}
Normalise the weights: $\omega_1=1/(1+3)=0.25$, $\omega_2=3/(1+3)=0.75$. Combine:
\begin{align*}
c &= 0.25\times0.20+0.75\times0.20 = 0.05+0.15 = 0.20.
\end{align*}
\end{solution}

\begin{solution}{ex:material:3}
\begin{align*}
\frac{\alpha}{F} &= \frac{0.02}{15} = 0.0013333,\\
p^* &= 1-0.0013333 = 0.9986667.
\end{align*}
By the union bound an ordinary pixel gets some man-made name with probability at most
$15\times0.0013333=0.02=\alpha$. Adding families raised the bar so that $\alpha$ keeps its meaning.
\end{solution}

\begin{solution}{ex:material:4}
The pixel centre is at column $j=20.5$, row $i=10.5$. With \cref{eq:material:affine}:
\begin{align*}
X &= 30\times20.5+0\times10.5+500{,}000 = 615+0+500{,}000 = 500{,}615\text{ m},\\
Y &= 0\times20.5+(-30)\times10.5+1{,}400{,}000 = 0-315+1{,}400{,}000 = 1{,}399{,}685\text{ m}.
\end{align*}
The identity transform would give $X=20.5$, $Y=10.5$: pixel coordinates, which a GIS would place near the
origin of the coordinate system, nowhere near the scene. That is why the CRS and transform must be recovered
at export.
\end{solution}

\begin{solution}{ex:material:5}
The action is accepted, but in the worker \code{_check_band_grid} finds that the vendable's wavelengths equal
the 165-band grid while V2 reads native bands, and the action fails with ``model 'Indradhanu V2' reads native
bands: run \code{band_filter_apply} with \code{use_common_wavelength_grid=false}''. Set
\code{use_common_wavelength_grid: false} on the band filter. MNF-RX, a classical detector, runs on either
grid, so it can stay in the same action. With today's default band filter (native bands, since 2026-10-31) the
action runs as it is.
\end{solution}

\begin{solution}{ex:material:6}
\begin{align*}
\text{total} &= 800\times900 = 720{,}000,\\
\text{valid} &= 720{,}000-20{,}000 = 700{,}000,\\
\texttt{excluded\_percent} &= \frac{100\times35{,}800}{700{,}000} = \frac{3{,}580{,}000}{700{,}000} = 5.114.
\end{align*}
Only cloud is excluded. Shadow is a class in the map but is not zeroed, so a vessel in cloud shadow is scored
normally by residual Mahalanobis; its score is not 0.
\end{solution}

\begin{solution}{ex:material:7}
\begin{align*}
P(\text{bare metal}) &= 0.7\times0.5=0.35, & P(\text{painted metal}) &= 0.7\times0.3=0.21,\\
P(\text{tarp}) &= 0.7\times0.2=0.14, & P(\text{natural}) &= 1-0.7=0.30,
\end{align*}
and $0.35+0.21+0.14+0.30=1.00$.
\end{solution}

\begin{solution}{ex:material:8}
$\vect{d}=1.0\times(12.5,-2,10)=(12.5,-2,10)$. Against soil: $\vect{s}\cdot\vect{d}=7.5\times12.5+3.75\times2+5\times10
=93.75+7.5+50=151.25$; $\vect{s}\cdot\vect{s}=95.3125$; $\vect{d}\cdot\vect{d}=260.25$; ACE
$=151.25^2/(95.3125\times260.25)=22876.56/24804.58=0.9223$. Without arithmetic: the fill only scales $\vect{d}$, and
ACE does not change when $\vect{d}$ is scaled. At fill 0.25: $0.7724\times0.7151=0.5523$, the table's 0.55 for level~2
right; it must agree, because a class can be right only when level~1 was.
\end{solution}

